
\documentclass{article}

%--- Document Formatting ---%
\usepackage{lmodern} % gives you the more modern Latin Modern fonts
\usepackage[T1]{fontenc} % fontenc provides better and more character encoding
\usepackage[margin=2.5cm, bottom=4cm, top=3cm, footskip=2cm]{geometry} % geometry allows changing page sizes, orientation and margin
\usepackage{microtype}
\usepackage[skip=8pt]{parskip} % better spacing between paragraphs
\usepackage{titling} % makes the title closer to the top
\setlength{\droptitle}{-5em}

\usepackage{titlesec} % titlesec allows changing the style, spacing and font of chapter, sections and subsections
% {left-indent}{space before}{space after}
\titlespacing*{\section}{0pt}{6ex plus 1ex minus .2ex}{2.3ex plus .2ex}
\titlespacing*{\subsection}{0pt}{2.5em plus 1ex minus .2ex}{2ex plus .2ex}
\linespread{1.15}

% Set distances between equations and text
\setlength{\jot}{8pt} % Distance between equations
\AtBeginDocument {
\setlength{\abovedisplayskip}{18pt}
\setlength{\belowdisplayskip}{18pt}
\setlength{\abovedisplayshortskip}{12pt}
\setlength{\belowdisplayshortskip}{12pt}
}

%--- Language ---%
\usepackage[australian]{babel}
\usepackage[useregional]{datetime2}


%--- Symbols ---%
\usepackage{amssymb} % amssymb provides symbols, mainly number systems
\usepackage{cancel} % adds \cancel{x}, allowing striking out derivations
%\usepackage{siunitx} % used for physics based units


%--- Tools and General Formatting ---%
\usepackage{mathtools} % an extension for amsmath, loads amsmath itself
\usepackage{amsthm} % improves proof based math
\usepackage{enumitem} % enumitem gives more control over lists

%\usepackage{graphicx} % use to include images/figures
%\usepackage{xcolor} % for colour
%\usepackage[hidelinks]{hyperref} % adds hyperlinks, can sometimes cause issues. Keep last
%\usepackage{cleveref} % extension for hyperref, auto formatted cross-reference



\begin{document}

\title{Parabolas}
\author{Montague Hamilton}
\date{\today}
\maketitle

\section{Questions}
\begin{enumerate}[itemsep=12pt]
  \item \textbf{Describe fully and test: $y = -3(x + 2)^{2} + 5$ - vertex, opening direction, shape and the y-intercept.}

    The vertex will occur at $(-2, 5)$, as
    \begin{align*}
      y &= -3((-2) + 2)^{2} + 5 \\
      y &= -3(0)^{2} + 5 \\
      y &= 5
    \end{align*}
    
    The opening direction will be downwards, as the y values get inverted by the negative in front of the expression, $-3(x+2)^{2}$.
    The shape will be quite narrow, as for every addtion to x, y changes by 3 (shrinks in this case).
    The y-intercept occurs when $x = 0$, so
    \begin{align*}
      y &= -3((0) + 2)^{2} + 5 \\
      y &= -3 \cdot 4 + 5 \\
      y &= -7
    \end{align*}
    The y-intercept occurs at $(0,-7)$.

  \item \textbf{Describe fully: $y = 2\sin(3x) -4$ - amplitude, midline and period.}

    The amplitude will be $2$, as for every value of $\sin(3x)$, y is twice as big as otherwise. The midline will be $-4$,
    as the last number added or subtracted from the rest of the function shifts every $y$ output by that amount. 
    The period will be $6\pi$, as the function will cycle $3$ times faster than otherwise. 

    \newpage
  \item \textbf{Spot the error: a student says $y = (x + 3)^{2}$ has vertex (3,0). 
    Test the claim, state the true vertex, and name the rule they violated.} 

    At $x = 3$,
    \begin{align*}
      y &= ((3) + 3)^{2} \\
      y &= 36
    \end{align*}
    At $x = 2$, $y = 25$ and at $x = 4$, $y = 49$, so it is clear that $x = 3$ can't be the vertex, as the $y$ values
    for $x = 2$ and $x = 4$ aren't the same. The vertex will be at $x = -3$, as
    \begin{align*}
      y &= ((-3) + 3)^{2} \\
      y &= (0)^{2}
    \end{align*}
    To check, $x = -4$, $y = 1$, should be equal to, $x = -2$, $y = 1$. So the vertex is at $(-3,0)$.
    The student most likely forgot the negative in front of the $3$.

  \item \textbf{Reverse engineer: a parabola has vertex $(2, -1)$ and passes through $(0,7)$.
    Find its equation in the form $y = a(x-h)^{2} +k$}

    If the vertex is at $(2, -1)$, then $h = 2$ and $k = -1$. To complete the equation we need to find $a$,
    plugging in the values for the known point $(0,7)$, 
    \begin{align*}
      7 &= a(0-2)^{2}-1 \\
      8 &= 4a \\
      a &= 2 
    \end{align*}
    So the final form of the equation would be,
    \[
      y = 2(x-2)^{2} - 1
    \]
    
  \item \textbf{Synthesis, one sentence each: describe the full journey from $y = x^{2}$ to $y = -(x-4)^{2} + 2$ -
    every knob in order, and where the vertex ends up.} 

    The $-4$ means the vertex starts at $x = 4$. The $2$ means the vertex starts at $y = 2$.
    The negative out the front of the parathesis means the parabola is inverted, but there isn't a number,
    so the graph isn't stretched or compressed. The vertex is at $(4,2)$.
    
    
\end{enumerate}


\end{document}
